\clearpage
\SourceLeaf{486}%
\fancyhead[L]{\small BERICHTE ÜBER DIE GRADMESSUNG.}
\pdfbookmark[0]{Bericht über die Resultate der trigonometrischen Messungen.}{bericht-resultate-1848}
\begin{center}\large [BERICHT ÜBER DIE RESULTATE DER TRIGONOMETRISCHEN MESSUNGEN.]\end{center}
\par
— — —

Von den Rechnungen, durch die der Übergang von dem rohen Messungs-Material [in den Journalen, welche 35 Hefte in der Reinschrift füllen] zu den Endresultaten [im Coordinaten-Verzeichniss] gemacht ist, habe ich nur einen dem Umfange nach sehr kleinen Theil unter meine jetzigen Vorlagen aufnehmen können.

In der That, wenn es möglich wäre, alle jene Rechnungen in extenso vollständig wieder aufzustellen, so möchten solche leicht vier oder sechsmal so viele Bände füllen [als jetzt die Journale binden]. Allein theils der Umstand, dass der grössere Theil der Details jener Rechnungen gar nicht aufbewahrt ist, theils die Form, in der sich die noch immer sehr voluminösen Fascikel der aufbewahrten Papiere befinden, haben zur Folge, dass eine vollständige und geordnete Wiederherstellung \textit{aller} Rechnungen fast dasselbe bedeuten würde, wie eine nochmalige Wiederholung meiner ganzen Arbeit. Ich habe mich demnach auf die geordnete Extrahirung desjenigen Theils der Zwischenrechnungen beschränkt, der als der prägnanteste und nützlichste betrachtet werden muss, nemlich auf die tabellarische Zusammenstellung aller an den verschiedenen Beobachtungsplätzen festgelegten und orientirten Richtungswinkel, wobei eine Parallele mit dem Göttinger Meridian den Nullpunkt oder Ausgangspunkt bildet. Diese tabellarischen Darstellungen sind unter der Benennung von \textit{Abrissen}\footnote{Die ganze Anzahl wird etwas über fünftehalbhundert betragen.} in sechs Heften zusammengeordnet, wobei ich von dem.Professor \textsc{Goldschmidt} mehrfache Beihülfe erhalten habe, welcher zugleich die Reinschriften grösstentheils selbst gemacht hat. — — —

Zur Entschuldigung der so sehr verspäteten Beendigung dieses Geschäfts muss ich noch bemerken, dass die Verspätung hauptsächlich daher entstanden ist, dass zur Erledigung der sowohl bei den Abschriften als noch mehr bei der Anfertigung der Abrisse aufstossenden zahlreichen Zweifel nicht selten erst langwierige Nachforschungen gemacht werden mussten. — — —

Sowie Zeit, Gesundheit und Kräfte es verstatten, werden meinen beiden ersten auf die geodätischen Probleme bezüglichen Abhandlungen noch ein Paar andere den speciellen Gegenständen noch näher tretend nachfolgen.

Göttingen den 15. März 1848.
\begin{flushright}C. F. \textsc{Gauss}.\end{flushright}
